% Reference count one at the write: the counted heap, the two calculi, the
% lowering strategy, the pair that characterises the target's write, and the
% typing that puts the count where a write needs it.
%
% Build, from this directory. Highlighting is minted's, so the build needs shell
% escape and Pygments on PATH:
%   python3 gen-laws.py     # the law frames and counts.tex, after a library change
%   python3 gen-frames.py   # the frames themselves: the quoted ranges, comment-free
%   latexmk -pdf -shell-escape lowering.tex
% Check that every frame still quotes the library, and that counts.tex holds the
% counts the library gives:
%   python3 check-listings.py
%
% pdflatex: the toolchain's lualatex cannot load its font loader, so the Unicode
% in the quoted Lean is mapped to font commands by newunicodechar below. A
% character in a quoted range with no mapping fails the build, which is the
% signal to add one.
\documentclass[11pt,a4paper]{article}

\usepackage[T1]{fontenc}
\usepackage[utf8]{inputenc}
\usepackage{lmodern}
\usepackage{amsmath,amssymb,amsthm,mathtools}
\usepackage{newunicodechar}
\newunicodechar{×}{\ensuremath{\times}}
\newunicodechar{·}{\ensuremath{\cdot}}
\newunicodechar{∘}{\ensuremath{\circ}}
\newunicodechar{α}{\ensuremath{\alpha}}
\newunicodechar{β}{\ensuremath{\beta}}
\newunicodechar{γ}{\ensuremath{\gamma}}
\newunicodechar{δ}{\ensuremath{\delta}}
\newunicodechar{ε}{\ensuremath{\varepsilon}}
\newunicodechar{ι}{\ensuremath{\iota}}
\newunicodechar{κ}{\ensuremath{\kappa}}
\newunicodechar{λ}{\ensuremath{\lambda}}
\newunicodechar{μ}{\ensuremath{\mu}}
\newunicodechar{ν}{\ensuremath{\nu}}
\newunicodechar{π}{\ensuremath{\pi}}
\newunicodechar{ρ}{\ensuremath{\rho}}
\newunicodechar{σ}{\ensuremath{\sigma}}
\newunicodechar{τ}{\ensuremath{\tau}}
\newunicodechar{φ}{\ensuremath{\varphi}}
\newunicodechar{χ}{\ensuremath{\chi}}
\newunicodechar{ω}{\ensuremath{\omega}}
\newunicodechar{Γ}{\ensuremath{\Gamma}}
\newunicodechar{Δ}{\ensuremath{\Delta}}
\newunicodechar{Σ}{\ensuremath{\Sigma}}
\newunicodechar{Φ}{\ensuremath{\Phi}}
\newunicodechar{ℓ}{\ensuremath{\ell}}
\newunicodechar{→}{\ensuremath{\to}}
\newunicodechar{←}{\ensuremath{\leftarrow}}
\newunicodechar{↔}{\ensuremath{\leftrightarrow}}
\newunicodechar{↦}{\ensuremath{\mapsto}}
\newunicodechar{⟶}{\ensuremath{\longrightarrow}}
\newunicodechar{∀}{\ensuremath{\forall}}
\newunicodechar{∃}{\ensuremath{\exists}}
\newunicodechar{∈}{\ensuremath{\in}}
\newunicodechar{∉}{\ensuremath{\notin}}
\newunicodechar{⊆}{\ensuremath{\subseteq}}
\newunicodechar{∧}{\ensuremath{\wedge}}
\newunicodechar{∨}{\ensuremath{\vee}}
\newunicodechar{¬}{\ensuremath{\neg}}
\newunicodechar{≠}{\ensuremath{\neq}}
\newunicodechar{≤}{\ensuremath{\leq}}
\newunicodechar{≥}{\ensuremath{\geq}}
\newunicodechar{⊥}{\ensuremath{\bot}}
\newunicodechar{∅}{\ensuremath{\emptyset}}
\newunicodechar{≃}{\ensuremath{\simeq}}
\newunicodechar{⟨}{\ensuremath{\langle}}
\newunicodechar{⟩}{\ensuremath{\rangle}}
\newunicodechar{…}{\ensuremath{\ldots}}
\newunicodechar{◾}{\ensuremath{\blacksquare}}
\newunicodechar{₀}{\textsubscript{0}}
\newunicodechar{₁}{\textsubscript{1}}
\newunicodechar{₂}{\textsubscript{2}}
\newunicodechar{′}{\textsuperscript{\prime}}
\usepackage{booktabs}
\usepackage{enumitem}
\usepackage[margin=1in]{geometry}
\usepackage{microtype}
\usepackage{xcolor}
\usepackage{minted}
% gruvbox-light: the most hues of the styles that stay legible on white paper.
\usemintedstyle{gruvbox-light}
\setminted{fontsize=\footnotesize,frame=single,rulecolor=\color{black!30},
           framesep=4pt,breaklines=true,baselinestretch=1.0}
% A quoted frame is written \codeframe[linerange=lo-hi]{../path}: a line range of
% a library file, with its comments removed, generated into frames/<key>.lean by
% gen-frames.py; frames/lookup.tex maps the range to that file. A call site
% therefore names the source it quotes and not the frame, and a range with no
% generated frame is a build error rather than a silently missing box. Each frame
% is one unbreakable box, so a frame that does not fit the space left on a page
% moves to the next one.
% Generated by gen-frames.py; do not edit by hand.
% Maps a \codeframe call's (path, range) to its generated frame, so a
% call site names the source and not the frame. The frames themselves are
% frames/<key>.lean, comment-free copies of the ranges, checked by
% check-listings.py against the library.

\expandafter\def\csname cf@../Scidonia/FracPerm/Ledger.lean@25-25\endcsname{9750611a7baa}
\expandafter\def\csname cf@../Scidonia/FracPerm/Ledger.lean@28-28\endcsname{5f928bda90e1}
\expandafter\def\csname cf@../Scidonia/FracPerm/Ledger.lean@34-35\endcsname{1913b63860bf}
\expandafter\def\csname cf@../Scidonia/FracPerm/StateVar/Value.lean@84-84\endcsname{dbe771f805da}
\expandafter\def\csname cf@../Scidonia/FracPerm/StateVar/Value.lean@96-97\endcsname{c011a33daef6}
\expandafter\def\csname cf@../Scidonia/FracPerm/StateVar/Value.lean@109-109\endcsname{c6b92c041b00}
\expandafter\def\csname cf@../Scidonia/FracPerm/StateVar/Value.lean@113-113\endcsname{9c6db65f9b82}
\expandafter\def\csname cf@../Scidonia/FracPerm/StateVar/Model.lean@369-370\endcsname{4ea30c66f461}
\expandafter\def\csname cf@../Scidonia/FracPerm/StateVar/Model.lean@423-424\endcsname{583a4eebbeaa}
\expandafter\def\csname cf@../Scidonia/FracPerm/StateVar/Model.lean@522-523\endcsname{bd6437b6df69}
\expandafter\def\csname cf@../Scidonia/FracPerm/StateVar/Model.lean@529-530\endcsname{a0fef9ee5433}
\expandafter\def\csname cf@../Scidonia/FracPerm/StateVar/Model.lean@535-535\endcsname{57d391554422}
\expandafter\def\csname cf@../Scidonia/FracPerm/StateVar/Model.lean@550-551\endcsname{b15f7df298db}
\expandafter\def\csname cf@../Scidonia/FracPerm/StateVar/Model.lean@597-600\endcsname{ef31681ee890}
\expandafter\def\csname cf@../Scidonia/FracPerm/StateVar/Model.lean@607-608\endcsname{e53944a7355a}
\expandafter\def\csname cf@../Scidonia/FracPerm/StateVar/Model.lean@582-587\endcsname{d4cc154e74a3}
\expandafter\def\csname cf@../Scidonia/FracPerm/StateVar/Model.lean@642-643\endcsname{14055dfba24e}
\expandafter\def\csname cf@../Scidonia/FracPerm/StateVar/Model.lean@656-657\endcsname{ba995450fed5}
\expandafter\def\csname cf@../Scidonia/FracPerm/StateVar/Model.lean@667-668\endcsname{611a3eae16f6}
\expandafter\def\csname cf@../Scidonia/FracPerm/StateVar/Heap.lean@102-103\endcsname{5b740228b3d1}
\expandafter\def\csname cf@../Scidonia/FracPerm/StateVar/Heap.lean@118-119\endcsname{9d55df76baec}
\expandafter\def\csname cf@../Scidonia/FracPerm/StateVar/Heap.lean@126-127\endcsname{548b3c12e0e5}
\expandafter\def\csname cf@../Scidonia/FracPerm/StateVar/Heap.lean@135-140\endcsname{f8af60b72496}
\expandafter\def\csname cf@../Scidonia/FracPerm/StateVar/Heap.lean@157-162\endcsname{8f1dc76ac072}
\expandafter\def\csname cf@../Scidonia/FracPerm/StateVar/Heap.lean@168-171\endcsname{ce7f92922125}
\expandafter\def\csname cf@../Scidonia/FracPerm/StateVar/Heap.lean@214-214\endcsname{c875ce15c57e}
\expandafter\def\csname cf@../Scidonia/FracPerm/StateVar/Array.lean@80-81\endcsname{5084cdad01e1}
\expandafter\def\csname cf@../Scidonia/FracPerm/StateVar/Array.lean@91-92\endcsname{486d61359463}
\expandafter\def\csname cf@../Scidonia/FracPerm/StateVar/Array.lean@113-114\endcsname{dda3851818c2}
\expandafter\def\csname cf@../Scidonia/FracPerm/StateVar/Array.lean@120-121\endcsname{8d41205432b3}
\expandafter\def\csname cf@../Scidonia/FracPerm/StateVar/Array.lean@165-166\endcsname{8041063aa6d3}
\expandafter\def\csname cf@../Scidonia/FracPerm/StateVar/Array.lean@194-195\endcsname{6a9ad9c69b16}
\expandafter\def\csname cf@../Scidonia/FracPerm/StateVar/Array.lean@216-220\endcsname{f98f33160e64}
\expandafter\def\csname cf@../Scidonia/FracPerm/StateVar/Universe.lean@86-87\endcsname{55bff5c33c9b}
\expandafter\def\csname cf@../Scidonia/FracPerm/StateVar/Universe.lean@103-103\endcsname{e23eb9c0c689}
\expandafter\def\csname cf@../Scidonia/FracPerm/StateVar/Universe.lean@153-154\endcsname{90e6c7c31d1e}
\expandafter\def\csname cf@../Scidonia/FracPerm/StateVar/Universe.lean@159-160\endcsname{6b6a3bfec516}
\expandafter\def\csname cf@../Scidonia/FracPerm/StateVar/Universe.lean@178-181\endcsname{e882d85a6188}
\expandafter\def\csname cf@../Scidonia/FracPerm/StateVar/Universe.lean@185-187\endcsname{e28b3a42614c}
\expandafter\def\csname cf@../Scidonia/FracPerm/StateVar/Universe.lean@193-199\endcsname{1fb0ee566945}
\expandafter\def\csname cf@../Scidonia/FracPerm/StateVar/Universe.lean@273-277\endcsname{960c08503b82}
\expandafter\def\csname cf@../Scidonia/FracPerm/StateVar/Slot.lean@67-67\endcsname{191b5865ef60}
\expandafter\def\csname cf@../Scidonia/FracPerm/StateVar/Slot.lean@73-73\endcsname{2e5a81e5133f}
\expandafter\def\csname cf@../Scidonia/FracPerm/StateVar/Slot.lean@82-82\endcsname{d66728bab821}
\expandafter\def\csname cf@../Scidonia/FracPerm/StateVar/Slot.lean@131-132\endcsname{777bd7ea6843}
\expandafter\def\csname cf@../Scidonia/FracPerm/StateVar/Slot.lean@152-156\endcsname{6f3e387a51d5}
\expandafter\def\csname cf@../Scidonia/FracPerm/StateVar/Slot.lean@169-174\endcsname{dca6a2e07dbf}
\expandafter\def\csname cf@../Scidonia/FracPerm/StateVar/Slot.lean@192-192\endcsname{8b3813d52365}
\expandafter\def\csname cf@../Scidonia/FracPerm/StateVar/Layout5.lean@190-191\endcsname{531d56988e55}
\expandafter\def\csname cf@../Scidonia/FracPerm/StateVar/Layout5.lean@194-195\endcsname{81d23c37433c}
\expandafter\def\csname cf@../Scidonia/FracPerm/StateVar/Layout5.lean@198-199\endcsname{3d4fd270215b}
\expandafter\def\csname cf@../Scidonia/FracPerm/StateVar/Layout5.lean@210-211\endcsname{3b81713f1639}
\expandafter\def\csname cf@../Scidonia/FracPerm/StateVar/Layout5.lean@221-222\endcsname{9381cefb1a86}
\expandafter\def\csname cf@../Scidonia/FracPerm/StateVar/Layout5.lean@270-271\endcsname{12f178d518f8}
\expandafter\def\csname cf@../Scidonia/FracPerm/StateVar/Layout5.lean@474-475\endcsname{2e37ef6f1c71}
\expandafter\def\csname cf@../Scidonia/FracPerm/StateVar/Layout5.lean@511-512\endcsname{851dd834f567}
\expandafter\def\csname cf@../Scidonia/FracPerm/StateVar/Layout5.lean@1027-1028\endcsname{0b9db9893a57}
\expandafter\def\csname cf@../Scidonia/FracPerm/StateVar/Layout5.lean@1055-1056\endcsname{79dee75ddae1}
\expandafter\def\csname cf@../Scidonia/FracPerm/StateVar/Layout5.lean@1243-1246\endcsname{c68bb8fcb258}
\expandafter\def\csname cf@../Scidonia/FracPerm/StateVar/Layout5.lean@1265-1267\endcsname{75d40ce95c38}
\expandafter\def\csname cf@../Scidonia/FracPerm/StateVar/Layout5.lean@1283-1284\endcsname{030cb2e59e4d}
\expandafter\def\csname cf@../Scidonia/FracPerm/StateVar/Layout5.lean@1361-1364\endcsname{0f90aede78e5}
\expandafter\def\csname cf@../Scidonia/FracPerm/StateVar/Layout5.lean@1379-1382\endcsname{7501cb075a09}
\expandafter\def\csname cf@../Scidonia/FracPerm/StateVar/Layout5.lean@1402-1405\endcsname{7f8964e9f0e6}
\expandafter\def\csname cf@../Scidonia/FracPerm/StateVar/Sealed.lean@125-127\endcsname{a150e108c8bc}
\expandafter\def\csname cf@../Scidonia/FracPerm/StateVar/Sealed.lean@132-135\endcsname{fa3821e2108b}
\expandafter\def\csname cf@../Scidonia/FracPerm/StateVar/Sealed.lean@152-153\endcsname{232c9c093968}
\expandafter\def\csname cf@../Scidonia/FracPerm/StateVar/Sealed.lean@157-158\endcsname{50d3a390ba9d}
\expandafter\def\csname cf@../Scidonia/FracPerm/StateVar/Sealed.lean@191-193\endcsname{865f8eda760b}
\expandafter\def\csname cf@../Scidonia/FracPerm/StateVar/Sealed.lean@199-202\endcsname{0af04bcb3c10}
\expandafter\def\csname cf@../Scidonia/FracPerm/StateVar/Sealed.lean@213-217\endcsname{4b2ed661cd73}
\expandafter\def\csname cf@../Scidonia/FracPerm/StateVar/Sealed.lean@234-237\endcsname{9b83be61cae0}
\expandafter\def\csname cf@../Scidonia/FracPerm/StateVar/Sealed.lean@251-253\endcsname{e1e96f07c3d9}
\expandafter\def\csname cf@../Scidonia/FracPerm/StateVar/RefCount.lean@740-741\endcsname{c2bd52ed1e21}
\expandafter\def\csname cf@../Scidonia/FracPerm/StateVar/RefCount.lean@765-765\endcsname{c1c6e144a166}
\expandafter\def\csname cf@../Scidonia/FracPerm/StateVar/RefCount.lean@771-771\endcsname{d9b5c6bdb52b}
\expandafter\def\csname cf@../Scidonia/FracPerm/StateVar/RefCount.lean@779-779\endcsname{fbe0489cba29}
\expandafter\def\csname cf@../Scidonia/FracPerm/StateVar/RefCount.lean@787-788\endcsname{163c10511c27}
\expandafter\def\csname cf@../Scidonia/FracPerm/StateVar/RefCount.lean@804-805\endcsname{aa9dea32f52b}
\expandafter\def\csname cf@../Scidonia/FracPerm/StateVar/RefCount.lean@819-823\endcsname{3352aac86792}
\expandafter\def\csname cf@../Scidonia/FracPerm/Structured/Model.lean@292-293\endcsname{331b6c9d7cb0}
\expandafter\def\csname cf@../Scidonia/FracPerm/Structured/Model.lean@297-298\endcsname{51d76a5567c4}
\expandafter\def\csname cf@../Scidonia/FracPerm/Structured/Model.lean@203-204\endcsname{b6db30baf018}
\expandafter\def\csname cf@../Scidonia/FracPerm/Structured/Model.lean@365-366\endcsname{8d3f6d606d28}
\expandafter\def\csname cf@../Scidonia/FracPerm/Structured/Model.lean@372-373\endcsname{0df1d9c5942f}
\expandafter\def\csname cf@../Scidonia/FracPerm/Structured/Model.lean@415-415\endcsname{4210bba08da7}
\expandafter\def\csname cf@../Scidonia/FracPerm/Structured/Model.lean@430-431\endcsname{469cb37dc3cf}
\expandafter\def\csname cf@../Scidonia/FracPerm/Structured/Model.lean@549-552\endcsname{5aa96b675f45}
\expandafter\def\csname cf@../Scidonia/FracPerm/Structured/Model.lean@559-560\endcsname{08b6b25ac8b0}
\expandafter\def\csname cf@../Scidonia/FracPerm/Structured/Model.lean@463-468\endcsname{4f499282495f}
\expandafter\def\csname cf@../Scidonia/FracPerm/Structured/Model.lean@565-567\endcsname{f03a2399da1a}
\expandafter\def\csname cf@../Scidonia/FracPerm/Structured/Model.lean@572-573\endcsname{3908df01cf88}
\expandafter\def\csname cf@../Scidonia/FracPerm/Structured/Model.lean@586-587\endcsname{9ff29af667c0}
\expandafter\def\csname cf@../ScidoniaTest/FracPerm.lean@96-96\endcsname{26c4a5c56142}
\expandafter\def\csname cf@../ScidoniaTest/FracPerm.lean@110-110\endcsname{c7072cc01773}
\expandafter\def\csname cf@../ScidoniaTest/FracPerm.lean@124-124\endcsname{667c4844b79c}
\expandafter\def\csname cf@../ScidoniaTest/FracPerm.lean@136-136\endcsname{af579ae4aafd}
\expandafter\def\csname cf@../ScidoniaTest/FracPerm.lean@149-149\endcsname{7cca48d0769a}
\expandafter\def\csname cf@../ScidoniaTest/FracPerm.lean@160-161\endcsname{58d6bb9503f6}
\expandafter\def\csname cf@../ScidoniaTest/FracPerm.lean@313-313\endcsname{84be31499770}
\expandafter\def\csname cf@../ScidoniaTest/StateVar.lean@281-281\endcsname{846203002cf7}
\expandafter\def\csname cf@../ScidoniaTest/StateVar.lean@287-287\endcsname{6247b6d47e02}
\expandafter\def\csname cf@../ScidoniaTest/StateVar.lean@293-293\endcsname{8ad7a3493477}
\expandafter\def\csname cf@../ScidoniaTest/Model5.lean@335-336\endcsname{3fb4624b4fe2}
\expandafter\def\csname cf@../ScidoniaTest/Model5.lean@353-353\endcsname{a10fbeeae676}
\expandafter\def\csname cf@../ScidoniaTest/Count.lean@44-44\endcsname{743326bf0f37}
\expandafter\def\csname cf@../ScidoniaTest/Count.lean@56-57\endcsname{ce08ec48a432}
\expandafter\def\csname cf@../ScidoniaTest/Count.lean@72-73\endcsname{6ef5df07669f}
\expandafter\def\csname cf@../ScidoniaTest/Count.lean@125-125\endcsname{a34eea8628ee}
\expandafter\def\csname cf@../ScidoniaTest/Count.lean@133-135\endcsname{018ed4c58639}
\expandafter\def\csname cf@../Scidonia/FracPerm/StateVar/Description.lean@24-46\endcsname{952632ccc0d9}
\expandafter\def\csname cf@../Scidonia/FracPerm/StateVar/Layout5.lean@91-119\endcsname{07c4092830a2}
\expandafter\def\csname cf@../Scidonia/FracPerm/StateVar/Layout5.lean@122-138\endcsname{4a24d4844560}
\expandafter\def\csname cf@../Scidonia/FracPerm/StateVar/Layout5.lean@147-166\endcsname{4f032786efd0}
\expandafter\def\csname cf@../Scidonia/FracPerm/StateVar/Layout5.lean@169-177\endcsname{0feaed59d5eb}
\expandafter\def\csname cf@../Scidonia/FracPerm/Ledger.lean@16-22\endcsname{4eb8cea25657}
\expandafter\def\csname cf@../Scidonia/FracPerm/Hoare.lean@107-132\endcsname{0303bf6cccb7}
\expandafter\def\csname cf@../Scidonia/FracPerm/Hoare.lean@145-152\endcsname{a45fb2b10d38}
\expandafter\def\csname cf@../Scidonia/FracPerm/Hoare.lean@178-184\endcsname{aef4c8295edc}
\expandafter\def\csname cf@../Scidonia/FracPerm/Hoare.lean@229-266\endcsname{8cc85285784b}
\expandafter\def\csname cf@../Scidonia/FracPerm/Counter.lean@22-49\endcsname{079f0ad3e32d}
\expandafter\def\csname cf@../Scidonia/FracPerm/StateVar/Value.lean@46-67\endcsname{bfb6d10b4c7b}
\expandafter\def\csname cf@../Scidonia/FracPerm/StateVar/Value.lean@68-82\endcsname{1f6a83f938ee}
\expandafter\def\csname cf@../Scidonia/FracPerm/StateVar/Class.lean@59-91\endcsname{3128a5f31a98}
\expandafter\def\csname cf@../Scidonia/FracPerm/StateVar/Heap.lean@44-47\endcsname{63b5d37b4f3c}
\expandafter\def\csname cf@../Scidonia/FracPerm/StateVar/Heap.lean@49-68\endcsname{0b25f74df449}
\expandafter\def\csname cf@../Scidonia/FracPerm/StateVar/Heap.lean@70-77\endcsname{6c7df423d965}
\expandafter\def\csname cf@../Scidonia/FracPerm/StateVar/Heap.lean@79-89\endcsname{1f6a8050c45e}
\expandafter\def\csname cf@../Scidonia/FracPerm/StateVar/Surface.lean@38-59\endcsname{2ff677e38d0a}
\expandafter\def\csname cf@../Scidonia/FracPerm/StateVar/Array.lean@57-76\endcsname{4d421a7cc316}
\expandafter\def\csname cf@../Scidonia/FracPerm/StateVar/Universe.lean@111-117\endcsname{55053fa59189}
\expandafter\def\csname cf@../Scidonia/FracPerm/StateVar/Universe.lean@210-240\endcsname{e1a2ef0a6af5}
\expandafter\def\csname cf@../Scidonia/FracPerm/StateVar/Slot.lean@46-63\endcsname{d3e2d380990e}
\expandafter\def\csname cf@../Scidonia/FracPerm/StateVar/Sealed.lean@52-60\endcsname{713c79572283}
\expandafter\def\csname cf@../Scidonia/FracPerm/StateVar/Sealed.lean@66-74\endcsname{9ec521f46635}
\expandafter\def\csname cf@../Scidonia/FracPerm/StateVar/Sealed.lean@83-90\endcsname{de0a942c6e6d}
\expandafter\def\csname cf@../Scidonia/FracPerm/StateVar/Layout5.lean@203-208\endcsname{85a1cf2fd656}
\expandafter\def\csname cf@../Scidonia/FracPerm/StateVar/Layout5.lean@298-308\endcsname{5af8e2fc8c99}
\expandafter\def\csname cf@../Scidonia/FracPerm/StateVar/Layout5.lean@457-471\endcsname{118d1cbd1fc1}
\expandafter\def\csname cf@../Scidonia/FracPerm/StateVar/Layout5.lean@539-555\endcsname{d010bd5a10fa}
\expandafter\def\csname cf@../Scidonia/FracPerm/StateVar/Layout5.lean@679-684\endcsname{61f9875d02fa}
\expandafter\def\csname cf@../Scidonia/FracPerm/StateVar/Layout5.lean@710-729\endcsname{95f71547697b}
\expandafter\def\csname cf@../Scidonia/FracPerm/StateVar/Layout5.lean@772-784\endcsname{20bf19049e08}
\expandafter\def\csname cf@../Scidonia/FracPerm/StateVar/RefCount.lean@713-728\endcsname{9729f794a150}
\expandafter\def\csname cf@../Scidonia/FracPerm/Structured/Surface.lean@127-149\endcsname{c9cdc256f152}
\expandafter\def\csname cf@../Scidonia/FracPerm/Structured/Surface.lean@166-179\endcsname{f8b42168ef1a}
\expandafter\def\csname cf@../Scidonia/FracPerm/Structured/Surface.lean@263-270\endcsname{dc0b62954321}

% counts.tex holds the numbers both papers in this directory state about the
% library and about their own panels. This paper's panel count is generated under
% its own name and read below under \panelcount, so the prose says "panel" and
% the numeral comes from the file; check-listings.py fails when that file is
% stale. The companion paper's counts keep their own names.
% Generated from the library by gen-laws.py; do not edit by hand.
% The numbers the paper's prose states about the library and about the
% frames it quotes, one \newcommand each, so no numeral is typed by hand.
% The paper's preamble reads this file, and check-listings.py recomputes
% it from the library and from the paper's own call sites.

\newcommand{\theoremcount}{694}  % public theorems under Scidonia/
\newcommand{\vendoredtheoremcount}{17}  % of them vendored from Lean's own model
\newcommand{\lawcount}{110}  % law statements the appendix reproduces
\newcommand{\panelcount}{37}  % panels the paper quotes
\newcommand{\appendixpanelcount}{32}  % of them in the appendix
\newcommand{\refcountframecount}{7}  % frames of the reference-count section

\let\panelcount\loweringpanelcount
\def\codeframe[linerange=#1-#2]#3{%
  \edef\cfkey{cf@#3@#1-#2}%
  \ifcsname\cfkey\endcsname\else
    \PackageError{paper}{no generated frame for #3:#1-#2}%
      {Run python3 gen-frames.py in Paper/.}\fi
  \edef\cffile{frames/\csname\cfkey\endcsname.lean}%
  \par\smallskip\noindent
  \begin{minipage}{\linewidth}%
    \inputminted{lean}{\cffile}%
  \end{minipage}\par}
% A path is set with \path so that it may break at its slashes.
\usepackage{url}\urlstyle{same}
% A frame is labelled with the file and lines it quotes:
\newcommand{\sourceline}[3]{\par\smallskip\noindent\textbf{#1}\hfill
  {\footnotesize\itshape\path{#2}, lines #3}\par\nobreak}
\usepackage[colorlinks=true,linkcolor=blue!45!black,citecolor=blue!45!black,urlcolor=blue!55!black]{hyperref}
\usepackage{cleveref}

\theoremstyle{definition}
\newtheorem{definition}{Definition}[section]
\newtheorem{remark}[definition]{Remark}
\theoremstyle{plain}
\newtheorem{theorem}[definition]{Theorem}
\newtheorem{lemma}[definition]{Lemma}

\newcommand{\lpure}{\textup{$\lambda$Pure}}
\newcommand{\lrc}{\textup{$\lambda$RC}}
\newcommand{\rc}{\mathrm{rc}}
\newcommand{\code}[1]{\texttt{\detokenize{#1}}}
\newcommand{\lref}[1]{\code{#1}}
\newcommand{\rulename}[1]{\textsc{#1}}

\title{Reference count one at the write\\[2pt]
\large The counted heap, the lowering, and the typing that puts the count there}
\author{Gavin Mendel-Gleason\\\small Scidonia}
\date{2026}

\begin{document}
\maketitle

\begin{abstract}
A reference-counted runtime's documented behaviour is to test its own count
before a write and update in place on its exclusive branch; that machine test is
distinct from the count-one predicate proved here for the counted heap.
Implementations decide this at run time, because one compiled body serves owning
and borrowing call sites, and the decision is taken per site, from what the
compiler can see there; the status of that compiler behaviour is stated in
\cref{sec:trusted}. This paper formalises a pure core, a reference-counted
target over a counted heap, and the lowering between them, and proves the pair
that characterises that target's write: in its counted heap, at count one the
store is destructive and keeps its location; above one it copies. The runtime's
own test of the count, and the correspondence between a model count and the
machine's, are stated below as assumptions and measurements with the limits they
hold under, and not as theorems. That pair reads the count of a \emph{reference}
in the term's own account of itself; the compiler-side account reads a count of
a \emph{name}, as the typing section explains (\cref{sec:typing}), with its
implementation link trusted as stated in \cref{sec:trusted}.
The second half of the paper closes the gap with a typing: the write's rule asks
the subject's name to be owned exactly once in the store's environment and not
borrowed, and the rules the term itself carries that bind that name ask it to be
used exactly once in the scope that binds it --- at an abstraction, at a binding,
and at each of an arm's binders --- and, at a binding and at each arm, to be
introduced for, or matched out of, a construction rather than another name. What
the typing establishes is conditional: it is a judgement on names, so it neither
makes a cell's reference count one nor rules out a holder of it outside the term,
and those are preconditions the correspondence and the take-out carry. The
programs a second holder of the cell would need --- the subject also used as the
written value, an arm binding it twice, a use of it in a sibling position, and a
binding of it to another name --- have no typing, each beside the same-shape
program that has one. A correspondence
between the typing's environment and the target's account reads the typing's
count of a name into the write theorem's count of a reference, the injectivity
that reading needs is derived rather than assumed where a run reaches it, and the
two routes meet at one step of an accepted run. Everything is mechanised in Lean~4
against the core library alone; this repository's own proofs take no step by
native decision and acquire no axiom, and the one vendored block whose native
evaluation does generate certificate axioms is named where it is used. The lowering
is proved total and deterministic, not correct; that is stated, and so is each
thing measured or trusted rather than proved.
\end{abstract}

\section{Introduction: the count at a write}
\label{sec:intro}

Lean's values are immutable. The term \code{a.set! i v} denotes a new array: same
length, same elements except at \code{i}. No term of the language can change a
value that already exists, and no program can tell whether some other array
happens to share a buffer with this one. The source language is pure, and the
compiler is free to make it fast.

The compiler is fast in a specific way. The runtime tests a cell's count before a
write and takes one of two branches, into the cell's own buffer or into a fresh
one, and that test and its branches are the runtime's documented behaviour rather
than anything proved here. What this paper proves is the same pair for a counted
heap it defines, and at the model's count: at one the store is destructive and
keeps its location, above one it copies. Everything this paper reads onto a
machine carries its status: the runtime's test is documented, the correspondence
between a model count of one and that test is a hypothesis no instance of which is
constructed, the shape a compiler emits for these write paths is checked against
one pinned toolchain's output rather than proved, and the in-place behaviour itself
is measured on compiled code. Each is named, with its limit, in \cref{sec:trusted}.
Both branches implement the same value, so nothing about correctness depends on
which one is chosen; what is at stake is allocation, and therefore speed.

\paragraph{The count, and where it comes from.} A reference-counted runtime keeps
one number per cell: the number of references to it. At a write to a cell, that
number is the sum of two things --- the references the cell's own slots hold, plus
the references held outside the heap by whatever is running. Those outside
references are, at the moment of the write, the ones the term itself holds
together with any handle the surrounding structure holds. That sum is a fact about
the term and its data, and the account of \cref{sec:account} is this paper's
statement of it. Whether a machine's number is that sum where a write runs is the
hypothesis \cref{sec:trusted} names rather than derives, and a compiler reaches the
question one site at a time.

\paragraph{What a compiler reads per site.} The compiler's own per-site
arrangement, whose status \cref{sec:trusted} names, is that one compiled body
serves an owning call site and a borrowing one, so the test is decided at each
site from what the compiler can see there: how many times a variable is used,
whether two names may denote one cell, and how an update is projected through a
record's fields. The test is exact when the analysis is exact, and conservative
otherwise. Nothing in this is specific to arrays: the same test guards any cell
of a counted heap.

\paragraph{The value.} An in-place update avoids an allocation and reuses the
buffer. It is not an optimisation of the meaning of the program --- the meaning is
the same either way --- it is an optimisation of the run.

\paragraph{The crux.} The condition a write would like to hold --- that the cell it
writes is the program's own, named nowhere else --- is a fact about the
\emph{compiled code} and not about the meaning of the program. A source term has
no way to state it: it is not a value, not a type, and not an equation. Getting it
to hold where a programmer wants a write to land in place is therefore a
\emph{discipline} on the source, and a discipline is what this paper builds after
the target is characterised: a term is typed only when its writes are expressible
on the target at all, and the typing is stated as a rule of the source so that the
programmer can see which terms qualify. The count the target reads is still a
count in a heap; the typing says which terms put the right count there.

\paragraph{What this paper does, and how each claim is labelled.} Sections
\ref{sec:pair} to \ref{sec:hypotheses} fix the mechanics: the counted heap, the
pure source calculus, the lowering strategy, the reference-counted target, the
account, and the pair of theorems that characterises the target's write. Section
\ref{sec:typing} is the typing, with its four refusals beside their typed
companions. Section \ref{sec:corr} reads the typing's count of a name into the
write theorem's count of a reference, and shows why the injectivity that reading
needs is necessary, and where it is derived rather than assumed. Section
\ref{sec:run} reads the same statement at a configuration an accepted run reaches.
Sections \ref{sec:trusted} and \ref{sec:future} state what is trusted and what is
not proved, and place the work.

Each claim below carries its status. \textbf{Proved} means a theorem of the
library, quoted here as a checked line range rather than transcribed;
\textbf{structural} means a fact the checked sources enforce without a single
theorem stating it; \textbf{measured} means an observation taken by running
compiled code; \textbf{trusted} means a link this development does not prove.

\section{The condition, and the pair}
\label{sec:pair}

This section fixes everything the pair of theorems in \cref{sec:pair-statement}
names. It is the mechanics the pair rests on: a counted heap with the two branches
of its write (\cref{sec:heap}), the source calculus the lowering starts from
(\cref{sec:source}), the lowering strategy (\cref{sec:lower}), the target's
runtime over the same heap (\cref{sec:target}), the account that relates a term to
the references it holds (\cref{sec:account}), and the two hypotheses the account
cannot supply (\cref{sec:hypotheses}).

\subsection{The pair}
\label{sec:pair-statement}

\noindent The write is destructive when its operand is owned once and nothing else
names the cell; otherwise it copies. That is the pair proved here, and it is the
only claim made about the target's write.

%!check: theorem store_owned_is_destructive {c : CHeap} {hs : List Ref} {ρ : TEnv} {m : Multi}
\sourceline{An owned store keeps its location}{Scidonia/RC/Ownership.lean}{559--564}
\codeframe[linerange=559-564]{../Scidonia/RC/Ownership.lean}

%!check: theorem store_shared_copies {c : CHeap} {hs : List Ref} {ρ : TEnv} {m : Multi} {y : Var}
\sourceline{A shared store copies}{Scidonia/RC/Ownership.lean}{600--606}
\codeframe[linerange=600-606]{../Scidonia/RC/Ownership.lean}

\noindent \lref{Couple} is the coupling invariant (\cref{sec:account}),
\lref{Owns} is the term-indexed account of the term's own references, \lref{expand
m} is the reference list the account stands for, and \lref{update} is the heap's
two-branch write (\cref{sec:heap}). The first theorem's conclusion is the
destructive branch read back: the location still holds the written slots. The
second's is the copying branch: the written slots are at a location the heap did
not list, the old cell keeps its slots, and the old cell loses a reference. The
two hypotheses the account cannot supply --- no slot and no run handle naming the
cell --- are named rather than hidden (\cref{sec:hypotheses}).

Read together the two theorems say: the write is in place exactly where the
account counts the operand once. The second is what keeps the first from being
vacuous, and \cref{sec:inj} exhibits the configuration that makes that precise: a
subject a second reference names \emph{must} copy, or the update would be visible
through the alias. A statement that every store is destructive is simply false,
and the pair is the honest form of the claim.

\subsection{The counted heap}
\label{sec:heap}

\noindent A cell holds the references its slots name and the count the runtime keeps
for it; a heap lists the live cells.

%!check: structure Cell where
\sourceline{A cell, and a heap}{Scidonia/FracPerm/StateVar/RCHeap.lean}{156--172}
\codeframe[linerange=156-172]{../Scidonia/FracPerm/StateVar/RCHeap.lean}

\noindent One reference names one cell, so the count a heap keeps at a reference is
not a choice among cells.

%!check: def CHeap.WF (c : CHeap) : Prop := (CHeap.refs c).Nodup
\sourceline{Well-formedness, and liveness}{Scidonia/FracPerm/StateVar/RCHeap.lean}{329--337}
\codeframe[linerange=329-337]{../Scidonia/FracPerm/StateVar/RCHeap.lean}

\noindent No live cell has count zero, because a release deallocates at zero. A
cell's count is what it must be: its occurrences among the live cells' slots, plus
the handles held outside the heap --- the references the running term holds.

%!check: def Counted (c : CHeap) (handles : List Ref) : Prop :=
\sourceline{What the counts mean}{Scidonia/FracPerm/StateVar/RCHeap.lean}{461--462}
\codeframe[linerange=461-462]{../Scidonia/FracPerm/StateVar/RCHeap.lean}

\noindent The write branches on that count: at one the cell keeps its location (and
the references its old slots held are dropped, cascading while any listed cell
reaches zero), otherwise a fresh cell takes the written slots and the old cell
loses a reference.

%!check: def update (c : CHeap) (handles : List Ref) (r : Ref) (ks : List Ref) : CHeap × List Ref :=
\sourceline{The heap's write, in two branches}{Scidonia/FracPerm/StateVar/RCHeap.lean}{1858--1863}
\codeframe[linerange=1858-1863]{../Scidonia/FracPerm/StateVar/RCHeap.lean}

\noindent The two branches, as already proved of the heap: a destructive write is
read back at the same location, a copying write at a location the heap did not list.

%!check: theorem update_preserves_location_of_exclusive {c : CHeap} {handles : List Ref}
\sourceline{The destructive branch}{Scidonia/FracPerm/StateVar/RCHeap.lean}{1945--1949}
\codeframe[linerange=1945-1949]{../Scidonia/FracPerm/StateVar/RCHeap.lean}

%!check: theorem update_copies_of_shared {c : CHeap} {handles : List Ref} (hc : Counted c handles)
\sourceline{The copying branch}{Scidonia/FracPerm/StateVar/RCHeap.lean}{1973--1979}
\codeframe[linerange=1973-1979]{../Scidonia/FracPerm/StateVar/RCHeap.lean}

\paragraph{Status.} Proved: both branches, at a heap that is well formed and at
handles the counts are stated for.

\subsection{The source calculus}
\label{sec:source}

\noindent One term type serves both calculi. Nodes are nullary, unary or binary; an
$n$-ary node is a spine. The write is a subject, a slot index and a value.

%!check: inductive Expr where
\sourceline{One term type for both calculi}{Scidonia/RC/Syntax.lean}{110--136}
\codeframe[linerange=110-136]{../Scidonia/RC/Syntax.lean}

\noindent The source, \lref{Pure}, is the fragment with neither instruction.

%!check: inductive Pure : Expr → Prop where
\sourceline{The source fragment}{Scidonia/RC/Syntax.lean}{140--159}
\codeframe[linerange=140-159]{../Scidonia/RC/Syntax.lean}

\noindent The free-variable function and the closedness discipline are what make
the substitution below sound: this is the discipline the source is used at.

%!check: def uses : Expr → List Var
\sourceline{Free occurrences, and closedness}{Scidonia/RC/Syntax.lean}{299--316}
\codeframe[linerange=299-316]{../Scidonia/RC/Syntax.lean}

\noindent Values are the canonical forms the reduction's rules match on.

%!check: inductive Value : Expr → Prop where
\sourceline{Values}{Scidonia/RC/Pure.lean}{151--161}
\codeframe[linerange=151-161]{../Scidonia/RC/Pure.lean}

\noindent The redexes: application of an abstraction, the two binders, a match on a
value node, and the write read as a functional update.

%!check: inductive Root : Expr → Expr → Prop where
\sourceline{The redexes}{Scidonia/RC/Pure.lean}{698--710}
\codeframe[linerange=698-710]{../Scidonia/RC/Pure.lean}

\noindent The reduction is a redex, or a step in one of the contexts below. There
is no rule inside a match's subject, a write's subject or a node's fields: the
source is written so that those are values already.

%!check: inductive Step : Expr → Expr → Prop where
\sourceline{The reduction}{Scidonia/RC/Pure.lean}{716--725}
\codeframe[linerange=716-725]{../Scidonia/RC/Pure.lean}

%!check: inductive Ectx where
\sourceline{The contexts the source steps in}{Scidonia/RC/Pure.lean}{837--845}
\codeframe[linerange=837-845]{../Scidonia/RC/Pure.lean}

\noindent Substitution is naive, and therefore capture free only on closed
arguments.

%!check: theorem uses_subst_of_closed : ∀ (e : Expr) (x : Var) (v : Expr), Pure e → Closed v →
\sourceline{Substitution on closed arguments}{Scidonia/RC/Pure.lean}{538--539}
\codeframe[linerange=538-539]{../Scidonia/RC/Pure.lean}

\noindent The reduction is a function, which is what lets the lowering be one too.

%!check: theorem step_deterministic : ∀ {e e₁ : Expr}, Step e e₁ →
\sourceline{The contractum is unique}{Scidonia/RC/Pure.lean}{781--782}
\codeframe[linerange=781-782]{../Scidonia/RC/Pure.lean}

\paragraph{Status.} Proved: determinism, and that a step preserves closedness
(\lref{step_preserves_closed}, beside the substitution law above).

\subsection{The lowering}
\label{sec:lower}

\noindent A derivation relates a source term to its target under a borrowed and an
owned multiset of occurrences. Availability is about names, not multiplicities: a
name used twice may be fitting at one owned occurrence, and the strategy supplies
the second use with a duplication.

%!check: def Fits (Δ Γ : Occs) (e : Expr) : Prop := ∀ x ∈ uses e, x ∈ Γ ∨ x ∈ Δ
\sourceline{Availability is about names}{Scidonia/RC/Lower.lean}{158}
\codeframe[linerange=158-158]{../Scidonia/RC/Lower.lean}

\noindent The relation: one rule per source construct. The splits are the functions
\lref{claimEnv}, \lref{bindEnv} and \lref{armEnv} rather than an existential, and
each rule's conclusion wraps its translation in \lref{wrap}, the prologue that
duplicates and releases.

%!check: inductive Lower : Occs → Occs → Expr → Expr → Prop where
\sourceline{The lowering, one rule per construct}{Scidonia/RC/Lower.lean}{232--272}
\codeframe[linerange=232-272]{../Scidonia/RC/Lower.lean}

\noindent The relation is exactly the graph of the computable strategy
\lref{lowerB} --- \lref{lowerB_sound} and \lref{lowerB_complete} are its two
directions --- so the strategy's own properties are the relation's. It is total on
the fragment and deterministic, and the second is why the splits are canonical:
with an arbitrary split of a binding's environment the one required duplication
could be placed in either subterm and the targets would differ.

%!check: theorem lower_total {Δ Γ : Occs} {e : Expr} (hp : Pure e) (hfit : Fits Δ Γ e) :
\sourceline{Totality}{Scidonia/RC/Lower.lean}{514--515}
\codeframe[linerange=514-515]{../Scidonia/RC/Lower.lean}

%!check: theorem lower_deterministic {Δ Γ : Occs} {e e₁ e₂ : Expr} (h₁ : Lower Δ Γ e e₁)
\sourceline{Determinism}{Scidonia/RC/Lower.lean}{524--526}
\codeframe[linerange=524-526]{../Scidonia/RC/Lower.lean}

\paragraph{Status.} Proved. These are properties of the translation. Nothing here
says the target computes what the source computes; that is
\cref{sec:trusted}.

\subsection{The target and its runtime}
\label{sec:target}

\noindent A configuration carries the heap, the handles the run holds, an
environment from variables to references, the term's own account of its
references, and the term. The environment is a list whose first pair for a name is
the one the name names.

%!check: abbrev TEnv := List (Var × Ref)
\sourceline{The runtime environment}{Scidonia/RC/Target.lean}{102}
\codeframe[linerange=102-102]{../Scidonia/RC/Target.lean}

\noindent The term's own account is one pair per occurrence the term itself holds,
and \lref{expand} projects the references it stands for --- the tags are not read
by the counts.

%!check: abbrev Multi := List (Var × Ref)
\sourceline{The term's own account}{Scidonia/RC/Target.lean}{178--182}
\codeframe[linerange=178-182]{../Scidonia/RC/Target.lean}

\noindent Each instruction moves the account in step with the heap: a duplication
increments the count and adds an occurrence, a release decrements, or deallocates
at a count of one.

%!check: inductive TStep :
\sourceline{The count operations, as steps}{Scidonia/RC/Target.lean}{541--558}
\codeframe[linerange=541-558]{../Scidonia/RC/Target.lean}

\noindent The write is the heap's own \lref{update}, given the account as outside
handles in addition to the run's. Its two rules are the two branches of
\lref{update}, and the only difference between them is the count.

%!check: | storeExclusive {c : CHeap} {hs : List Ref} {ρ : TEnv} {m : Multi} {y : Var} {i : Nat}
\sourceline{The two store rules}{Scidonia/RC/Target.lean}{603--625}
\codeframe[linerange=603-625]{../Scidonia/RC/Target.lean}

\noindent Each of the two reads the subject's occurrence out of the term's own
account (\lref{hsub}), so an account that does not hold it cannot step. The
position of the write in the tree arrives through the context, not through the
typing: the positions the runtime steps in are the contexts below.

%!check: inductive TCtx where
\sourceline{The positions the target steps in}{Scidonia/RC/Target.lean}{671--685}
\codeframe[linerange=671-685]{../Scidonia/RC/Target.lean}

%!check: theorem tstep_of_plug : ∀ {E : TCtx} {c : CHeap} {hs : List Ref} {ρ : TEnv} {m : Multi}
\sourceline{A step inside a well-formed context is a step}{Scidonia/RC/Target.lean}{719--722}
\codeframe[linerange=719-722]{../Scidonia/RC/Target.lean}

\paragraph{Status.} Proved: that the congruence rules are these contexts one by
one, which \lref{tstep_of_plug} states for a well-formed context.

\subsection{The account}
\label{sec:account}

\noindent Exact tagging: each pair names the reference its variable is bound to,
not merely some reference. The weaker reading admits a mis-tagged pair and makes
the theorems below arithmetic.

%!check: def Bound (ρ : TEnv) (m : Multi) : Prop := ∀ p ∈ m, ρ.lookup p.1 = some p.2
\sourceline{Exact tagging}{Scidonia/RC/Ownership.lean}{116}
\codeframe[linerange=116-116]{../Scidonia/RC/Ownership.lean}

\noindent The account is of the term: each consuming instruction reads its
occurrence out of it, and a store reads the occurrence its subject names.

%!check: inductive Owns (ρ : TEnv) : Multi → Expr → Prop where
\sourceline{The account is of the term}{Scidonia/RC/Target.lean}{417--435}
\codeframe[linerange=417-435]{../Scidonia/RC/Target.lean}

\noindent The coupling invariant: the term's own references are handles. That is
the whole content of the middle conjunct.

%!check: def Couple (c : CHeap) (hs : List Ref) (ρ : TEnv) (m : Multi) : Prop :=
\sourceline{The coupling invariant}{Scidonia/RC/Ownership.lean}{161--162}
\codeframe[linerange=161-162]{../Scidonia/RC/Ownership.lean}

\noindent Every step preserves it, contexts included.

%!check: theorem couple_step {c : CHeap} {hs : List Ref} {ρ : TEnv} {m : Multi} {e : Expr}
\sourceline{The invariant is preserved}{Scidonia/RC/Ownership.lean}{489--491}
\codeframe[linerange=489-491]{../Scidonia/RC/Ownership.lean}

\paragraph{Status.} Proved, for the relation's steps.

\subsection{The two hypotheses, and where they come from}
\label{sec:hypotheses}

\noindent \lref{store_owned_is_destructive} cannot derive, from the account, that
no heap slot and no run handle names the cell. Both are what a nested update's
take-out provides: at a write the run resolves a cell, takes the element out of
the slot the write lands at, holds it as a handle, and checks that cell before
recursing --- so the addressed cell has no slot naming it, and the run's own
handle is the only name on it.

%!check: theorem exclusive_head {c : CHeap} {hs : List Ref} {r : Ref}
\sourceline{The focus check is what the pair names}{Scidonia/FracPerm/StateVar/Run.lean}{1019--1021}
\codeframe[linerange=1019-1021]{../Scidonia/FracPerm/StateVar/Run.lean}

\noindent The take-out also pairs the extracted occurrence with a handle and drops
that handle before the focus check, which is what leaves the cell exclusive again
at the count of one the write needs.

%!check: theorem takeout_dropped_count_one {c : CHeap} {e : Ref} (hc : Counted c [e, e])
\sourceline{The take-out's own arithmetic}{Scidonia/FracPerm/StateVar/RCHeap.lean}{2082--2085}
\codeframe[linerange=2082-2085]{../Scidonia/FracPerm/StateVar/RCHeap.lean}

\noindent The mechanisation names the two surviving facts among the theorem's
hypotheses rather than deriving them, and cites the take-out lemmas that supply
them. \cref{sec:run} reads the same two facts off a write step of an accepted run.

\section{The typing}
\label{sec:typing}

The pair of \cref{sec:pair-statement} counts a \emph{reference} in the term's own
account. A compiler counting at a site counts the occurrences of a \emph{name} in
the environment that site is compiled in. The two agree only because the target's
accounts happen to be injective on references, and nothing in the source language
says they are. What the source needs, and what this section builds, is a judgement
that admits a write only where the environment gives its subject exactly one
occurrence, and that refuses the shapes where a second holder of the cell exists.

\noindent The judgement is the lowering's own environment discipline read as a
typing: the same clauses over the same canonical splits, one rule per source
construct, no target term, the write's clause strengthened to what a destructive
write needs, and two further conditions on the name a write's subject carries.

\subsection{The write's clause}
\label{sec:clause}

\noindent The write's rule is the only place where multiplicity and not membership
is read. It requires the subject's variable to be owned exactly once in the
environment of the store and not borrowed.

%!check: | var {Δ Γ : Occs} {x : Var} (h : x ∈ Γ ∨ x ∈ Δ) : Typed Δ Γ (.var x)
\sourceline{A variable is typed where an environment supplies it}{Scidonia/RC/Typing.lean}{255--264}
\codeframe[linerange=255-264]{../Scidonia/RC/Typing.lean}

%!check: | store {Δ Γ : Occs} {y : Var} {i : Nat} {v : Expr} (hown : Γ.count y = 1)
\sourceline{The write's clause}{Scidonia/RC/Typing.lean}{299--301}
\codeframe[linerange=299-301]{../Scidonia/RC/Typing.lean}

\noindent \lref{Γ.count y = 1} is a count of the subject's \emph{name} in the
store's environment, and \lref{y ∉ Δ} says the name is not borrowed there. Three
readings of one judgement follow from it, and the first is the one the theorems
above are read through: every write in a typed term has, at the environment the
typing assigns its position, a subject owned exactly once there and not borrowed.
That is the traversal \lref{WriteOwned}, whose recursion is the typing's own.

%!check: def WriteOwned (Δ Γ : Occs) : Expr → Prop
\sourceline{Every write's subject is owned once, carried by a traversal}{Scidonia/RC/Typing.lean}{313--331}
\codeframe[linerange=313-331]{../Scidonia/RC/Typing.lean}

\paragraph{Status.} Proved. \lref{typed_write_owned} carries the traversal over a
typed term, and \lref{typed_lowers} shows a typed term is one the lowering strategy
answers from the same environments: the typing is a restriction of the discipline
the lowering walks, and not a second discipline beside it.

\subsection{The two conditions a written name carries}
\label{sec:names}

\noindent The clause above reads the environment \emph{at the store}, and an
environment is a filter of its parent's, so a name's occurrences are claimed
separately by each position that uses them: shares overlap, and a sibling's use of
the subject's name is not counted at the store. A term whose write's subject is
also used in a sibling position therefore meets the write's clause --- the store's
own occurrence is the subject's --- while the program as a whole uses the subject's
name twice, and the lowering supplies the sibling's use with a duplication. Read at
the store's environment alone, the typing cannot see the program's multiplicity, so
two further conditions are needed, and both are side conditions of the rules that
\emph{bind} the name.

\noindent The first condition --- a written name is used exactly once in the scope
that binds it --- is carried by each rule of the term that binds a name, and the
abstraction carries it alone.

%!check: | lam {Δ Γ : Occs} {x : Var} {body : Expr}
\sourceline{The first condition alone, on an abstraction}{Scidonia/RC/Typing.lean}{262--264}
\codeframe[linerange=262-264]{../Scidonia/RC/Typing.lean}

\noindent The abstraction cannot carry the second: its binder is bound to the
argument an application supplies, and no side condition of an abstraction or of an
application can see that argument's shape, so there is no term of the rule to ask
whether it is a construction (\cref{sec:reach}). A binding carries both, for the one
name it binds.

%!check: | bind {Δ Γ : Occs} {x : Var} {e₁ e₂ : Expr}
\sourceline{Both conditions, on a binding}{Scidonia/RC/Typing.lean}{273--277}
\codeframe[linerange=273-277]{../Scidonia/RC/Typing.lean}

\noindent An arm carries both for each of its binders, and the construction the
name is matched out of is the matched subject.

%!check: | matchC {Δ Γ : Occs} {s : Expr} {c : Ctor} {ys : List Var} {body : Expr}
\sourceline{Both conditions, on an arm}{Scidonia/RC/Typing.lean}{282--287}
\codeframe[linerange=282-287]{../Scidonia/RC/Typing.lean}

\noindent The first condition is what makes the \emph{program} use the bound name
once in the scope that binds it, rather than counting only the store's share of the
environment; \lref{writes} is the traversal the condition reads, true at a write
whose subject is the free variable \lref{x} and stopped by a binder of \lref{x}.
The second says the written name comes from a construction rather than another
name: a binding's written name is bound to a node the program built and never to a
variable, which would make the written name a second name beside that variable
whatever cell it names. \lref{typed_write_linear} and \lref{typed_write_built}
carry the two conditions over a typed term, as traversals in the style of
\lref{WriteOwned}.

\paragraph{Status.} Proved: the conditions are side conditions of the rules that
bind the name, and are carried over typed terms, so the write's own clause and the
two name conditions are three readings of one judgement.

\subsection{The refusals, and their typed companions}
\label{sec:refusals}

\noindent Each shape in which a second holder of the cell exists has no typing, and
each sits beside the same-shape program that has one. The eight programs are these.

%!check: def siblingUse : Expr := .lam 0 (.con1 7 (.con2 8 (.store (.var 0) 0 (.con0 9)) (.var 0)))
\sourceline{The eight programs}{Scidonia/RC/Typing.lean}{541--572}
\codeframe[linerange=541-572]{../Scidonia/RC/Typing.lean}

\paragraph{The subject is also the written value.} The store's environment carries
two occurrences of the subject's name, and the write's clause asks for one.

%!check: theorem dupValue_no_typing : ¬ Typed [] [] dupValue := by
\sourceline{The subject used twice: no typing}{Scidonia/RC/Typing.lean}{577}
\codeframe[linerange=577-577]{../Scidonia/RC/Typing.lean}

%!check: theorem ownedValue_typed : Typed [] [] ownedValue := by
\sourceline{The same shape, the subject used once}{Scidonia/RC/Typing.lean}{597}
\codeframe[linerange=597-597]{../Scidonia/RC/Typing.lean}

\paragraph{An arm binds the subject twice.} The subject is supplied twice by the
binder list, so the body uses the bound name twice.

%!check: theorem dupArm_no_typing : ¬ Typed [] [] dupArm := by
\sourceline{The arm binding it twice: no typing}{Scidonia/RC/Typing.lean}{588}
\codeframe[linerange=588-588]{../Scidonia/RC/Typing.lean}

%!check: theorem ownedArm_typed : Typed [] [] ownedArm := by
\sourceline{The same arm, binding it once}{Scidonia/RC/Typing.lean}{605}
\codeframe[linerange=605-605]{../Scidonia/RC/Typing.lean}

\paragraph{A sibling use.} The store's own clause is met --- the sibling's use is
not counted at the store --- and the program still uses the subject's name twice,
so the first name condition refuses it. This is the refutation that the store's own
environment alone cannot give, and its companion writes through the name and does
not use it elsewhere.

%!check: theorem sibling_use_no_typing : ¬ Typed [] [] siblingUse := by
\sourceline{A sibling use: no typing}{Scidonia/RC/Typing.lean}{626}
\codeframe[linerange=626-626]{../Scidonia/RC/Typing.lean}

%!check: theorem sibling_use_owned : Typed [] [] siblingOwned := by
\sourceline{The same program, using the name once}{Scidonia/RC/Typing.lean}{642}
\codeframe[linerange=642-642]{../Scidonia/RC/Typing.lean}

\paragraph{An alias binding.} A binding's written name is bound to another name, so
the written name is a second name beside it whatever cell it names. The refutation
is at every environment, because the condition is a statement about the term, and
its companion binds the name to a node the program built.

%!check: theorem alias_bind_no_typing {Δ Γ : Occs} {x z : Var} {i : Nat} {v : Expr} :
\sourceline{A binding to another name: no typing}{Scidonia/RC/Typing.lean}{655--656}
\codeframe[linerange=655-656]{../Scidonia/RC/Typing.lean}

%!check: theorem built_bind_typed : Typed [] [] builtBind := by
\sourceline{The same binding, to a construction}{Scidonia/RC/Typing.lean}{667}
\codeframe[linerange=667-667]{../Scidonia/RC/Typing.lean}

\paragraph{The reading.} The shapes where a second holder of the cell exists have
\emph{no} typing, and each beside its companion shows why: the companions differ
from the refuted programs in nothing but the multiplicity of the subject's name or
the shape of what binds it. Where the term itself binds the name a write's subject
carries, the rule that binds it uses that name exactly once in the scope that binds
it and, where the rule is a binding or an arm, introduces it for a construction
rather than another name: the name's multiplicity in that scope is one, and what
binds it is not another name for the same cell. That is the source-side half of the
condition the target's destructive branch asks the account for. It is not a count on
the target heap, and it is not a statement that no other holder of the cell exists:
the judgement is on names, a store is typed in an environment whose two names may
name one reference while \lref{Couple} holds (\cref{sec:inj}), and a subject that is
a free variable, an abstraction's argument or a field a binder names is a holder the
term's own rules cannot see (\cref{sec:reach}). The correspondence carries those
cases as its own preconditions, and the take-out supplies them where a run reaches
them (\cref{sec:corr}). \lref{dup_no_typing} and
\lref{drop_no_typing} record the other half: a resource instruction has no rule, so
a typed term mentions neither.

\subsection{What the conditions do not reach}
\label{sec:reach}

\noindent The two conditions are stated of the rules that bind a name \emph{within}
the term, and a written name may be one the term binds nowhere. As the module says
of its own limitation: a subject that is a free variable of the term is the
environment's, and the store rule admits it wherever the environment holds it once
and it is not borrowed; an abstraction's name is bound to the argument an
application supplies, and no side condition of an abstraction or an application can
see that argument's shape; and a matching arm's binder names a \emph{field} of the
matched cell, and a field that is itself a name is a second name for that name's
cell. Those cases --- together with the write reached through a nested update, where
the cell's count of one and the absence of any other name for it are facts the
take-out provides and \lref{store_owned_is_destructive} names among its hypotheses
rather than deriving --- remain outside the judgement.

What is proved is the name side of the write: where the term itself binds the name
a subject carries, that name is used exactly once in the scope that binds it and is
introduced for, or matched out of, a construction.

\paragraph{Status.} Proved, of the rules the term itself carries, and explicitly
not of the three cases above.

\subsection{The typing is a statement about the source}
\label{sec:at-source}

\noindent A target-side sufficient condition --- a judgement over the target's own
relation, asking at a store for the two hypotheses the write theorem needs --- is
the obvious other route, and it is not the one taken. It is refuted by the
lowering's own prologue: \lref{Runs} demands of a binding whose bound term is a
value that the continuation run at every reference the binder could take, and at
the well-built reference the account counts that reference three times.

%!check: theorem runs_store_owned {ρ : TEnv} {m : Multi} {y : Var} {i : Nat} {v : Expr}
\sourceline{A target-side condition, and what it costs}{Scidonia/RC/Discipline.lean}{403--405}
\codeframe[linerange=403-405]{../Scidonia/RC/Discipline.lean}

\noindent A source term the strategy walks and its canonical lowering are one pair
of witnesses to that:

%!check: def source : Expr := .bind 0 (.con2 7 (.var 1) (.var 2)) (.store (.var 0) 0 (.var 2))
\sourceline{A source term and its canonical lowering}{ScidoniaTest/Target.lean}{505--510}
\codeframe[linerange=505-510]{../ScidoniaTest/Target.lean}

%!check: theorem source_typed : Typed [] [0, 1, 2] source := by
\sourceline{The source is typed}{ScidoniaTest/Target.lean}{597}
\codeframe[linerange=597-597]{../ScidoniaTest/Target.lean}

%!check: theorem lowered_not_runs : ¬ Runs env0 [(0, 0), (1, 1), (2, 2)] lowered := by
\sourceline{The lowering is not judged}{ScidoniaTest/Target.lean}{666}
\codeframe[linerange=666-666]{../ScidoniaTest/Target.lean}

\noindent The typing holds at the environments the strategy walks, and the
target-side judgement does not hold at the lowering --- which is why the typing is
stated at the source, where the rules that bind a written name are visible, rather
than at the target, where the prologue's duplication has already hidden the
program's multiplicity behind an account.

\section{From names to the count}
\label{sec:corr}

The typing counts names in an environment; the write theorem of
\cref{sec:pair-statement} counts a reference in the term's own account. The link
between them is a correspondence between the target's environment \lref{ρ}, the
term's own account \lref{m} and the typing's environment \lref{Γ}.

\subsection{The correspondence}
\label{sec:bridge}

\noindent The correspondence has three clauses: every pair of the account names
the reference the environment binds its name to, the account's multiplicity at each
name is the typing's count of that name, and the account's tags are injective on
references.

%!check: structure Corr (ρ : TEnv) (m : Multi) (Γ : Occs) : Prop where
\sourceline{The correspondence}{Scidonia/RC/Correspondence.lean}{129--136}
\codeframe[linerange=129-136]{../Scidonia/RC/Correspondence.lean}

\noindent The first two clauses are already enough for half of the write theorem's
premises: the account's occurrence of the subject's name and the environment's
reading of it give the term-indexed account \lref{Owns}. This half needs no
injectivity, and it is why the countermodel of \cref{sec:inj} separates on the
count of the reference and not on the occurrence.

%!check: theorem owns_of_correspondence {ρ : TEnv} {m : Multi} {Γ : Occs} {y : Var} {r : Ref}
\sourceline{Half the bridge, without injectivity}{Scidonia/RC/Correspondence.lean}{292--295}
\codeframe[linerange=292-295]{../Scidonia/RC/Correspondence.lean}

\noindent The third clause is what turns the typing's count of a \emph{name} into
the write theorem's count of a \emph{reference}: the account has one pair for the
name, and an injective account puts no second name on that pair's reference.

%!check: theorem bridge {ρ : TEnv} {m : Multi} {Γ : Occs} {y : Var} {r : Ref} {i : Nat} {v : Expr}
\sourceline{The bridge}{Scidonia/RC/Correspondence.lean}{305--307}
\codeframe[linerange=305-307]{../Scidonia/RC/Correspondence.lean}

\noindent Composed with the write theorem, a store typed at \lref{Γ} whose
environment counts its subject once, at a configuration corresponding in the sense
of \lref{Corr}, is destructive.

%!check: theorem store_destructive_of_typed {c : CHeap} {hs : List Ref} {ρ : TEnv} {m : Multi}
\sourceline{A typed store is destructive}{Scidonia/RC/Correspondence.lean}{435--440}
\codeframe[linerange=435-440]{../Scidonia/RC/Correspondence.lean}

\noindent Its hypotheses \lref{hslot}, \lref{hrun} and \lref{hks} are the write
theorem's own, named there as what the discipline's take-out provides and what the
write rule's check demands; the conclusion is that theorem's. \lref{bridge} reads
the typing's own write clause off the derivation (\lref{Typed.store_count}), so no
injectivity is an input of this theorem's either --- it is supplied by \lref{Corr},
and \cref{sec:takeout} is where that clause comes from at a run.

\paragraph{Status.} Proved.

\subsection{The injectivity is necessary}
\label{sec:inj}

\noindent The injectivity clause is the one neither the typing nor the coupling
supplies, and the configuration that separates them is two names naming one
reference: the store is typed there, the account is \lref{Bound} and inside
\lref{Couple}, the account's multiplicity at the subject's name is the typing's
count of it, and the reference is still counted twice. So the count of one does not
follow from the judgement and the invariant together.

%!check: theorem couple_does_not_give_count_one :
\sourceline{The whole invariant, and the count of two}{Scidonia/RC/Correspondence.lean}{404--408}
\codeframe[linerange=404-408]{../Scidonia/RC/Correspondence.lean}

\noindent Only \lref{RefInj} refuses that configuration, and the module exhibits
the refusal beside the three facts that do not (\lref{alias_typed},
\lref{alias_couple}, \lref{alias_not_inj}, with the reference count of two read at
\lref{alias_ref_count_two}). What would supply the clause on the environment's side
is a distinctness or freshness condition on the environment the program starts
from; the next subsection is where the run supplies it instead.

\paragraph{Status.} Proved: the countermodel is a theorem, and so is the failure
of the count of one at it.

\subsection{The injectivity at a take-out}
\label{sec:takeout}

\noindent At a nested write the injectivity is not an input. The run's counting at
the take-out's level, beside the coupling's own counting of one heap, forces the
account's references to be the one occurrence the take-out extracted.

%!check: theorem expand_eq_singleton_of_counted {c : CHeap} {hs : List Ref} {m : Multi} {e : Ref}
\sourceline{Two countings force the account's references}{Scidonia/RC/Transfer.lean}{233--235}
\codeframe[linerange=233-235]{../Scidonia/RC/Transfer.lean}

\noindent And an account whose references are a single occurrence is injective
whatever names its pairs carry.

%!check: theorem refInj_of_counted {c : CHeap} {hs : List Ref} {m : Multi} {e : Ref}
\sourceline{And one occurrence is injective}{Scidonia/RC/Transfer.lean}{264--265}
\codeframe[linerange=264-265]{../Scidonia/RC/Transfer.lean}

\noindent So at a take-out the injectivity is \emph{derived}, not assumed: the
take-out pins the account's size, and a one-pair account cannot name a reference
twice. The boundary is sharp, and the library states it too: two accounts with the
same references are one configuration to the run's side --- the same count at every
heap and handle list --- and one of them names a reference twice.

%!check: theorem refInj_not_determined :
\sourceline{The run's side does not determine the injectivity}{Scidonia/RC/Transfer.lean}{338--340}
\codeframe[linerange=338-340]{../Scidonia/RC/Transfer.lean}

\paragraph{Status.} Proved. The reading is exact in both directions: injectivity is
necessary (\cref{sec:inj}) and it is not a function of the references alone.

\section{Reading a run through it}
\label{sec:run}

The composition above is a statement about a configuration, and a run is a sequence
of them. This section reads the two routes at one step: the typing's route to the
write theorem, and the run model's own write step.

\noindent At a nested write the take-out has emptied the slot the write lands at and
holds the extracted occurrence as a handle, so the run's handle list is that
occurrence followed by the outer handles while the coupling's split is the outer
handles followed by the account's references. The two read the same references
exactly when the account is the extracted occurrence, and that is the coupling the
take-out hands the nested write.

%!check: theorem couple_of_takeout {c : CHeap} {hs : List Ref} {ρ : TEnv} {m : Multi} {r : Ref}
\sourceline{The coupling the take-out hands the nested write}{Scidonia/RC/Transfer.lean}{415--420}
\codeframe[linerange=415-420]{../Scidonia/RC/Transfer.lean}

\noindent The account is read \emph{out of} the run's handle list rather than added
beside it, which is why the coupling's handles are the run's own
(\lref{couple_at_run_handles_is_empty} at the coupling's split). The account's
identity is that theorem's hypothesis, since it is the coupling being built; where
an already-built coupling is at hand the identity is derived instead, from that
coupling's own counting beside the run's (\cref{sec:takeout}).

%!check: theorem writeStep_nested_destructive_of_couple {root : Ref} {c : CHeap} {hs : List Ref}
\sourceline{A write step of an accepted run, read through the coupling}{Scidonia/RC/Transfer.lean}{507--519}
\codeframe[linerange=507-519]{../Scidonia/RC/Transfer.lean}

\noindent A write step of an accepted configuration, its take-out, the account
through the coupling, and the schema's own checks are the write theorem's two
premises, so the destructive branch at the nested subject is a statement about that
step, with \lref{Run.writeStep_spec}'s own result beside it. The correspondence's
injectivity clause is built from the two countings rather than carried
(\lref{store_destructive_of_corr}). At a write whose path's level is its last, the
recursion's entry is the write's own subject and the leaf's check supplies the write
rule's premise on the written slots, so the two routes' conclusions are about one
cell.

%!check: theorem writeStep_subject_destructive_of_couple {root : Ref} {c : CHeap} {hs : List Ref}
\sourceline{The two routes at one cell}{Scidonia/RC/Transfer.lean}{540--552}
\codeframe[linerange=540-552]{../Scidonia/RC/Transfer.lean}

\paragraph{What this does, and what it does not.} It links the typing's route to the
run's route at one step of an accepted run. The induction over a whole program's
steps is the run model's own (\lref{Run.every_write_step_keeps_its_address}, the
per-run statement that a write's address holds the written slots afterwards, and
\lref{Run.writeStep_lands}, its per-step form), and neither of those is read through
the coupling. What is not built is that induction \emph{carrying the coupling}: the
run model's configurations carry no term, so the account the coupling relates to is
the caller's, and its term-side facts --- that the account's pairs name what the
environment binds them to, and that the account is the extracted occurrence --- are
supplied at the step rather than carried through the run.

\section{What is trusted, and what is not proved}
\label{sec:trusted}

\noindent \textbf{What is trusted.} Three links, each with an exact boundary.

\emph{The take-out's contribution.} That a nested update's cell has no slot and no
handle naming it is a property of the nesting, and is taken here as the two named
hypotheses of \lref{store_owned_is_destructive}; \cref{sec:hypotheses} is where the
run supplies them, and \cref{sec:run} is where they are read off a step.

\emph{The compiler's lowering.} That a compiler turns a language's write paths into
the shape modelled here is not proved, and the link is \emph{trusted}. The
accompanying repository bounds it with a check rather than a theorem: the check
reads the post-reference-counting intermediate representation the pinned toolchain
emits and asserts the shape on one executing path, failing the build when it does
not hold. That is evidence about the output of one pinned toolchain, at the
declarations the check names and the one path it selects; it is not a theorem about
compilers, and it does not discharge the link, which stays an assumption here.

\emph{The runtime's primitive.} The branch that stores in place is a runtime test
whose meaning is documented in the runtime and not modelled here. The counted model
states the predicate a write tests.

%!check: def IsExclusive (rc : Int32) : Prop := isSt rc ∧ rc = 1
\sourceline{The predicate a write tests, at a count}{Scidonia/FracPerm/StateVar/RefCount.lean}{720}
\codeframe[linerange=720-720]{../Scidonia/FracPerm/StateVar/RefCount.lean}

\noindent The model's count of one is not the machine's test in either direction,
and the difference is stated where it matters: the count's module says of the
store's run-time primitive that its effect on the count is an \emph{assumption} at
the one place the development needs it, labelled as such, and its write-path
theorem carries exactly those assumptions as hypotheses.

%!check: theorem takeOut_refCount_eq_one (rc0 rc1 rc2 : Int32)
\sourceline{The write path's hypotheses are assumptions}{Scidonia/FracPerm/StateVar/RefCount.lean}{819--823}
\codeframe[linerange=819-823]{../Scidonia/FracPerm/StateVar/RefCount.lean}

\noindent \emph{The in-place check.} The check a client can turn on is instrumented
and measured, not proved. In the logic it is the identity, so no specification
mentions it; its compiled implementation reports and exits when the object it is
given is not exclusively owned.

%!check: @[implemented_by assertExclusiveImpl]
\sourceline{The in-place check: the identity in the logic}{Scidonia/FracPerm/InPlace.lean}{92--93}
\codeframe[linerange=92-93]{../Scidonia/FracPerm/InPlace.lean}

\noindent It is measured on compiled code by the repository's own harness, which
runs a program under the check and reports the sites where an update would have
copied. That an update copies at such a site is a fact about the run that was
observed, not a theorem about the program.

\paragraph{The access-shape limit.} The store of \cref{sec:source} addresses one
index into a cell, and depth arrives through the term: the target's write rules are
read at the cell a variable names, and a nested write's path is a context over the
relation rather than a form of the term. A nested write's uniqueness is therefore
supplied by the take-out (\cref{sec:hypotheses}) and not by the typing, which is the
honest scope of \cref{sec:typing}: the typing speaks of the writes the term
addresses directly.

\noindent \textbf{What is not proved.} The lowering's semantic correctness: no
theorem says the target computes what the source computes, or that it is garbage
free. What is proved of the lowering is totality and determinism. The declarative
relation of the literature is not the one formalised here: the strategy's splits
are canonical, for determinism, so the results are properties of \emph{this}
strategy, which is what a compiler implements. And nothing here is about a
compiled program: the references are names in a list, not addresses.

\section{Mechanisation, and related work}
\label{sec:future}

\noindent The library publishes \theoremcount{} public theorems under
\code{Scidonia/}, of which \vendoredtheoremcount{} are vendored from Lean's own
model of the count. This paper quotes \panelcount{} panels of it, and every one is
a line range of the source rather than a transcription: the range is generated into
a frame with its comments removed, a \texttt{\%!check:} directive names a fragment
the range must contain, and \code{check-listings.py} fails when either has drifted.
The counts the prose states are read from a generated file, so a count that moves
with the library is regenerated rather than corrected by hand.

The development is a few thousand lines of Lean~4 against the core library only,
with no mathematics library dependency. This repository's own proofs have no
\code{sorry} and take no step by native decision, and every representative theorem
of this repository's has its axiom list asserted, so a proof of this repository's
that acquired a foreign axiom fails the build. The counts module is the exception,
and it is upstream's text: it is vendored from Lean's own model of the reference
count. Thirteen of its seventeen public theorems carry at least one generated
certificate axiom of the form \code{<theorem>._native.bv_decide.ax_*} beside
\code{propext}, \code{Classical.choice} and \code{Quot.sound}, because
\code{bv_decide} decides by native evaluation and asserts the result as an axiom;
four print none, and of those only \code{LEAN_RC_STUCK_ST_eq} invokes
\code{bv_decide} at all, its goal having been closed before the certificate step,
the other three being \code{grind} and \code{omega} proofs. Those axioms are not
this repository's to remove, and the vendored arithmetic is left as it stands. A
fact derived from a vendored theorem inherits that theorem's axioms and says which
declaration it inherits them from, and the checks of \code{ScidoniaTest/Count.lean}
pin the list \texttt{\#print axioms} emits for this repository's own facts, so any
other axiom fails the build there. Two design facts the formalisation forced are
recorded as toolchain facts: terms are a flat inductive, because a nested one
cannot be inducted on or derive decidable equality; and an evaluated check needs a
structural recursion, because the kernel does not reduce a well-founded one.

\noindent \textbf{Reference counting with reuse.} In-place update of a shared value,
and the precise counting with ownership that makes it possible, is the design of
Lean~4's runtime \cite{ullrich19counting} and is formalised as the linear resource
calculus $\lambda^1$ in the Pe\TeX{} paper \cite{reinking21perceus}, which proves
the counting sound and garbage free. This development follows that shape and
differs where this paper's subject differs: the deliverable is the count at the
write, as a pair, with the hypotheses the account cannot supply named; a typing that
puts the count where a write needs it, with the shapes a second holder would need
refuted; and the translation is a canonical strategy rather than the declarative
relation.

\noindent \textbf{Linearity, and formal models of counting.} An ordinary type system
does not track uniqueness: Linear Haskell \cite{bernardy18linear} gives a linear
arrow, and its guarantee is about the use of a value, not about a count in a heap.
A formal model of reference counting for code generation appears in F\'erey and
Shankar \cite{ferrey16codegen}. Neither is a count-one guarantee for a write, and
neither is stated as a condition on the name a write's subject carries.

\noindent \textbf{Future work.} Three steps: the lowering's semantic correctness
(result agreement and garbage freedom); carrying the coupling through a run's steps,
which needs the run model's configurations to carry the term, so that the two
routes of \cref{sec:run} compose over a whole program rather than at one step; and
a verified account, or a soundness theorem for the check, of the compiler's pass.

\begin{thebibliography}{9}

\bibitem{ullrich19counting}
Sebastian Ullrich and Leonardo de Moura.
\newblock Counting immutable beans: reference counting optimized for purely
functional programming.
\newblock In \emph{Proceedings of the 31st Symposium on Implementation and
Application of Functional Languages (IFL '19)}, 2019. arXiv:1908.05647.

\bibitem{reinking21perceus}
Alex Reinking, Ningning Xie, Leonardo de Moura, and Daan Leijen.
\newblock Pe\TeX{}: garbage free reference counting with reuse.
\newblock In \emph{Proceedings of the 42nd ACM SIGPLAN International Conference on
Programming Language Design and Implementation (PLDI '21)}, 2021.
\newblock Extended version: Microsoft Research Technical Report MSR-TR-2020-42,
which contains the proofs.

\bibitem{bernardy18linear}
Jean-Philippe Bernardy, Mathieu Boespflug, Ryan R. Newton, Simon Peyton Jones, and
Arnaud Spiwack.
\newblock Linear Haskell: practical linearity in a higher-order polymorphic
language.
\newblock In \emph{Proceedings of the 45th ACM SIGPLAN Symposium on Principles of
Programming Languages (POPL '18)}, 2018.

\bibitem{ferrey16codegen}
Gaspard F\'erey and Natarajan Shankar.
\newblock Code generation using a formal model of reference counting.
\newblock In \emph{NASA Formal Methods (NFM 2016)}, pages 150--165, 2016.

\bibitem{demoura21lean4}
Leonardo de Moura and Sebastian Ullrich.
\newblock The Lean 4 theorem prover and programming language.
\newblock In \emph{Automated Deduction (CADE-28)}, 2021.

\bibitem{scidonia-lean}
Gavin Mendel-Gleason.
\newblock The mechanisation of this paper, in the Scidonia \code{lean} repository:
\code{Scidonia/RC/}, the counted heap in
\code{Scidonia/FracPerm/StateVar/RCHeap.lean}, the run model in
\code{Scidonia/FracPerm/StateVar/Run.lean}, the check in
\code{Scidonia/FracPerm/InPlace.lean}, and the checks in
\code{ScidoniaTest/Typing.lean}, \code{ScidoniaTest/Correspondence.lean},
\code{ScidoniaTest/Transfer.lean} and \code{ScidoniaTest/Target.lean}.
\newblock \url{https://github.com/scidonia/lean-examples}.

\end{thebibliography}

\end{document}
